\documentclass[a4paper, 12pt]{article}
\usepackage{amsmath, amsfonts, amsthm, amssymb}
\usepackage{color, fullpage, hyperref}
\usepackage{graphicx}

\setlength{\parindent}{0pt}
\begin{document}
\textbf{\underline{Definition}} Group

Group is a set equipped with binary operations (usually denoted by 
$\cdot$ or $+$). The operation must satisfy the
\textit{group axioms}:
\begin{enumerate}
    \item Closure: $\forall a,b \in G, a\cdot b \in G$
    \item Associativity: $\forall a,b,c \in G, a\cdot (b\cdot c) = 
        (a\cdot b)\cdot c$
    \item Identity: $\forall a\in G,\exists e\in G, s.t~  
        e\cdot a = a\cdot e = a$
    \item Inverse: $\forall a\in G,\exists a^{-1}\in G, s.t~  
        a^{-1}\cdot a = a\cdot a^{-1}= a$
    \item Uniqueness: the identity and inverse element must be unique.
\end{enumerate}

\textbf{\underline{Definition}} Order of Group

The \textit{order of group}, denoted by $|G|$, is the number of element in $G$.

\vspace{5mm}

\textbf{\underline{Definition}} Order of Element

The \textit{order of element}, denoted by $|g|$, where $g\in G$, is the smallest
$n$ such that $g^n = e$.

\vspace{5mm}

\textbf{\underline{Definition}} Center of Group

An element $g\in G$ is in the \textit{center} of the group, $Z(G)$, iff $g$ and
commute with all elements in $G$. Formally:
$$G := \{g\in G|~ \forall a\in G, ag = ga\}$$

\vspace{5mm}

\textbf{\underline{Definition}} Centralizer

Let $a\in G$. The \textit{centralizer} of $a$ is defined as:

$$C(a) := \{g\in G|~ ag = ga\}$$


\vspace{5mm}

\textbf{\underline{Corollary}}

$G$ is abelian if and only if $Z(G) = G$.

\vspace{5mm}

\textbf{\underline{Definition}}

\vspace{5mm}

\textbf{\underline{Equivalence Relation}}

Given a set $S$. An \textit{equivalence relation} is a set $R\subset S \times S$
, such that $R = \{(a, b)|~ a,b\in S\}$, that satisfies the following:
\begin{enumerate}
    \item Reflexivity $\forall a\in S, (a, a)\in R$
    \item Symmetry $\forall a,b\in S, (a, b)\in R \implies (b, a)\in R$
    \item Transitivity $\forall a, b, c\in S, (a, b), (b, c)\in R \implies 
        (a, c)\in R$
\end{enumerate}

In general, $(a, b)\in R$ is written as $a\sim b$.
\newpage

\textbf{\underline{Definition}} Equivalent Class

The \textit{equivalence class} of $a\in G$, is the set of all element that 
equivalent to $a$.

$$[a] = \{g\in G|~ a\sim G\}$$

Given an equivalence relation on the set $S$, we can form equivalence classes in
the set $S$
\begin{enumerate}
    \item Each element $a$ belongs to an equivalence class, which is equivalence 
        class of $a$ itself.

        $\forall a\in G,~ a\in[a]$
    \item The equivalent classes are disjoint.

        $\forall a, b\in G,$ either $[a] = [b]$ or $[a]\cap [b] = \emptyset$

    \item The equivalence classes of each element partition the whole group

        $$\bigcup _{a\in X} [a] = X$$
\end{enumerate}

\textbf{\underline{Definition}} Dihedral Group

The usual \textit{representation} of $D_n$ is given by:
$$D_n := \langle r,s | r^n = e, s^2 = e, sr^k = r^{-k}s \rangle$$

We can interpret $r^j$ and $sr^j$ as rotation $\frac{360}{n}$ degree clockwise
and reflection along different axis. However, that is not the intrinsic property
for dihedral groups.

\vspace{5mm}
\textbf{\underline{Proposition}}

If $n$ is even, $Z(D_n) = \{e, r^{n/2}\}$, otherwise $Z(D_n) = \{e\}$.

\vspace{5mm}

\textbf{\underline{Definition}} Subgroup
Given a group $G$, $H$ is the \textit{subgroup} of $G$, written as $H < G$ iff
\begin{enumerate}
    \item $H\subset G$
    \item $e\in H$
    \item $\forall a, b\in H, ab\in H$
    \item $\forall a\in H, \exists!a^{-1}\in H$
\end{enumerate}


\end{document}
